%% 02-background.tex --- included by ../sfc.tex
\section{Background}
\label{sec:background}

\subsection{Quantum Circuits and Routing}
\label{sec:qcircuits}

A circuit acts on $n$ logical qubits through a sequence of gates; we
work in the restricted basis $\{\texttt{id}, \texttt{rz}, \texttt{sx},
\texttt{x}, \texttt{cx}\}$ of MQT-Bench's Qiskit-opt3-transpiled
circuits~\cite{quetschlich2023mqtbench}, i.e.\ arbitrary one-qubit
gates plus the two-qubit CX. A circuit is a directed acyclic graph
(DAG): for every logical qubit, consecutive gates acting on it are
joined by a dependency edge. The \emph{front} is the set of gates with
no unexecuted predecessor; only front gates can execute or be routed
towards execution.

Physical qubits are partitioned into \emph{cores}; each core has an
intra-core coupling graph, and pairs of cores may be joined by an
inter-core link whose two endpoints are designated \emph{communication
ports}. A \emph{layout} maps every logical qubit to a physical
position. A two-qubit gate is \emph{local} if both operands share a
core and \emph{remote} otherwise; routing a local gate that is not yet
adjacent requires intra-core SWAPs, and routing a remote gate requires
an inter-core communication primitive. We write $\kappa_\ell(q)$ for
the core holding $q$ under layout $\ell$, $\dL(c,p,p')$ for
intra-core distance and $\dC(c,c')$ for core-graph distance.

\subsection{Communication Primitives}
\label{sec:primitives}

Three inter-core primitives are in play, sketched in
Figure~\ref{fig:primitives}; each consumes exactly one EPR pair.

\begin{figure}[t]
\centering
% ---- (a) TeleData ----
\begin{tikzpicture}[x=0.42cm,y=0.42cm]
  % left core (2x3), port = right-middle; right core, port = left-middle
  \foreach \ox in {0,3.4}{
    \draw[rounded corners=2pt,black!50] (\ox-0.5,0.5) rectangle (\ox+1.5,-2.5);
    \foreach \xx in {0,1} \foreach \yy in {0,1,2}
      \node[slot] at (\ox+\xx,-\yy) {};
  }
  \node[portslot] at (1,-1) {};
  \node[portslot] at (3.4,-1) {};
  \draw[eprlink] (1,-1) -- (3.4,-1);
  \node[font=\tiny,red!70!black] at (2.2,-0.35) {EPR};
  \node[occ] at (0,-1) {};
  \node[font=\tiny,anchor=east] at (-0.35,-1) {$q$};
  \draw[-{Stealth[length=1.8mm]},semithick,black!70]
    (0,-0.75) .. controls (1.2,0.9) and (2.4,0.9) .. (3.4,-0.72);
  \node[font=\scriptsize] at (2.2,-3.3) {(a) \TeleData};
\end{tikzpicture}\hfill
% ---- (b) TeleGate ----
\begin{tikzpicture}[x=0.42cm,y=0.42cm]
  \foreach \ox in {0,3.4}{
    \draw[rounded corners=2pt,black!50] (\ox-0.5,0.5) rectangle (\ox+1.5,-2.5);
    \foreach \xx in {0,1} \foreach \yy in {0,1,2}
      \node[slot] at (\ox+\xx,-\yy) {};
  }
  \node[portslot] at (1,-1) {};
  \node[portslot] at (3.4,-1) {};
  \draw[eprlink] (1,-1) -- (3.4,-1);
  \node[font=\tiny,red!70!black] at (2.2,-0.35) {EPR};
  \node[occ] at (0,-1) {};
  \node[occ] at (4.4,-1) {};
  \node[font=\tiny,anchor=east] at (-0.35,-1) {$q_c$};
  \node[font=\tiny,anchor=west] at (4.75,-1) {$q_t$};
  \draw[catline] (0,-1) -- (1,-1);
  \draw[catline] (3.4,-1) -- (4.4,-1);
  \filldraw[blue!60!black] (0,-1) circle (1.4pt);
  \draw[blue!60!black] (4.4,-1) circle (2.6pt)
    +(0,-2.6pt) -- +(0,2.6pt) +(-2.6pt,0) -- +(2.6pt,0);
  \node[font=\scriptsize] at (2.2,-3.3) {(b) \TeleGate};
\end{tikzpicture}\hfill
% ---- (c) TeleGatePacket ----
\begin{tikzpicture}[x=0.42cm,y=0.42cm]
  \foreach \ox in {0,3.4}{
    \draw[rounded corners=2pt,black!50] (\ox-0.5,0.5) rectangle (\ox+1.5,-2.5);
    \foreach \xx in {0,1} \foreach \yy in {0,1,2}
      \node[slot] at (\ox+\xx,-\yy) {};
  }
  \node[portslot] at (1,-1) {};
  \node[portslot] at (3.4,-1) {};
  \draw[eprlink] (1,-1) -- (3.4,-1);
  \node[font=\tiny,red!70!black] at (2.2,-0.35) {EPR};
  \node[occ] at (0,-1) {};
  \node[font=\tiny,anchor=east] at (-0.35,-1) {$q_c$};
  \node[occ] at (4.4,0) {};
  \node[occ] at (4.4,-1) {};
  \node[occ] at (4.4,-2) {};
  \node[font=\tiny,anchor=west] at (4.75,0) {$t_1$};
  \node[font=\tiny,anchor=west] at (4.75,-1) {$t_2$};
  \node[font=\tiny,anchor=west] at (4.75,-2) {$t_3$};
  \draw[catline] (0,-1) -- (1,-1);
  \filldraw[blue!60!black] (0,-1) circle (1.4pt);
  \draw[catline] (3.4,-1) .. controls (3.9,-0.2) .. (4.4,0);
  \draw[catline] (3.4,-1) -- (4.4,-1);
  \draw[catline] (3.4,-1) .. controls (3.9,-1.8) .. (4.4,-2);
  \node[font=\scriptsize] at (2.2,-3.3) {(c) \TeleGatePacket};
\end{tikzpicture}
\caption{The three communication primitives; each consumes one EPR
  pair between the two communication ports (squares). (a)~\TeleData{}
  moves qubit $q$ (which must sit adjacent to its free port) onto the
  free destination port. (b)~\TeleGate{} executes one remote CX
  through a cat-entangler~\cite{yimsiriwattana2004generalized}
  (\appCatent{}): the target-side port acts as a proxy for
  the control $q_c$; neither qubit moves. (c)~\TeleGatePacket{}
  brackets an ordered run of remote CX gates sharing the control $q_c$
  between one cat-entangler/disentangler pair, amortising a single
  EPR pair over the whole run; each target is routed next to the port
  in turn while the ports stay reserved.}
\label{fig:primitives}
\end{figure}

\textbf{State teleportation (\TeleData{}).} Consumes one EPR pair to
move a logical qubit from a communication port in one core to a free
port in another, after which every later gate that qubit meets in its
new core is local. Its value is therefore \emph{latent}: it is paid
once and repaid over however many future gates the relocation makes
local.

\textbf{Gate teleportation (\TeleGate{}).} Executes a single remote CX
with neither operand moving. One EPR pair entangles the two
communication ports; a \emph{cat-entangler} ---
$\mathrm{CX}(q_c,e_A)$, measurement of $e_A$, and a classically
controlled $X$ on $e_B$ --- leaves $(q_c,e_B)$ in the cat state
$\alpha|00\rangle + \beta|11\rangle$, so every CX controlled by the
remote port $e_B$ acts as if controlled by $q_c$. A
\emph{cat-disentangler} --- Hadamard on $e_B$, measurement, a
classically controlled $Z$ on $q_c$ --- restores $q_c$ and frees both
ports (\appCatent{}). Its value is entirely \emph{direct}: exactly one
gate, and nothing is left behind.

\textbf{Same-control packets (\TeleGatePacket{}).} When several remote
CX gates share the same control and their targets lie in one core, the
run can be bracketed between a \emph{single} cat-entangler and
disentangler: the ports stay reserved, each target is routed beside
$e_B$ in turn, and the whole run costs one EPR pair instead of one
each. This is the amortisation idea of AutoComm~\cite{wu2022autocomm},
specialised to the shared-control case and applied inside a routing
loop rather than as a pre-pass.

The three differ in exactly one respect that matters for scoring: how
many gates one pair is spread over. \TeleData{} amortises over an
unknown future, \TeleGate{} over exactly one gate, and
\TeleGatePacket{} over a run whose length the compiler must discover.
Section~\ref{sec:sfc-scoring} makes that difference explicit as a
direct/latent/lookahead decomposition, and Section~\ref{sec:ablation}
measures whether the richer vocabulary is worth having.

\textbf{\SABRE{}.} \SABRE{}~\cite{li2019tackling} routes a single-core
circuit by repeatedly selecting the SWAP that minimises a front-layer
distance term plus a weighted extended-set lookahead term, both
expressed as $\Delta$ (old distance minus new distance) on gates sharing
a qubit with the candidate SWAP. \PORTER{} inherits this scoring form
for its intra-core stage (Section~\ref{sec:sfc-intra}) and extends the
same decomposition to its inter-core plan scorer
(Section~\ref{sec:sfc-scoring}).

%%────────────────────────────────────────────────────────────────────────────
